\section{Conclusion}

The paper asked three questions about the absolute-versus-relative choice.
Why does relativisation help for some KPIs and harm others?
How can a relative feature with $\eta<1$ still inform?
Does the same pairing geometry appear beyond sport?

Help and harm follow from two properties of each paired KPI: the variance ratio $\kappa=\sigma^{2}_{\mathrm{B}}/\sigma^{2}_{\mathrm{A}}$ and the within-match correlation $\rho$.
The Paired Efficiency Factor $\eta=(1+\kappa)/(1+\kappa-2\sqrt{\kappa}\,\rho)$ compares the variance of the difference with unpaired measurement.
The formula is distribution-free and contains Fisher's equal-variance result as the case $\kappa=1$.
Positive $\rho$ reduces that variance ($\eta>1$); negative $\rho$ inflates it ($\eta<1$).
That is why some KPIs should be relativised on efficiency grounds and others should not.
Efficiency is not the same as information about the outcome.
Under bivariate normality, $\eta$ and $I(X;Y)$ are related through the signal strength $\delta/\sigma_{\mathrm{A}}$, the standardised mean difference between the two entities.
A relative feature can therefore improve prediction when $\eta<1$, provided the mean separation remains large relative to the amplified noise.

For the rugby and football KPIs, $\kappa$ is the variance ratio of the two teams.
The correlation $\rho$ measures whether their values rise and fall together or against each other.
Those coordinates place every KPI in one of four quadrants.
Rugby and football occupy all four and concentrate in Quadrant~3, where $\kappa<1$ and $\rho<0$, so absolute features are preferred.
Four exemplars at comparable $\delta/\sigma_{\mathrm{A}}$ confirm both help and harm (\cref{tab:exemplars}).
The same $(\kappa,\rho)$ geometry appears in healthcare, clinical genomics, finance, and manufacturing, with different mean $\hat\eta$ (\cref{tab:validation}).
Quadrant membership is a first-pass guide for the feature choice.
Signal strength $\delta/\sigma_{\mathrm{A}}$ then scales how much of that geometry appears in prediction.
In Quadrant~4 it is required before a relative feature is retained.
